\documentclass{article}

\usepackage[utf8]{inputenc}
\usepackage[T1]{fontenc}
\usepackage{graphicx}
\usepackage{float}
\usepackage{amsmath}
\usepackage{amssymb}

\title{\textbf{Equação de Advecção-Difusão}\\
Modelação em Física e Engenharia}

\author{Tiago Bernardo - 53117 \and Gabriel Marques - 53087}

\date{Maio 2026}

\begin{document}

\maketitle

\section{Introdução}

A equação de advecção--difusão constitui uma das equações diferenciais parciais mais relevantes na modelação matemática de fenómenos físicos envolvendo transporte e dissipação.

Esta equação descreve a evolução temporal e espacial de uma grandeza conservada sujeita simultaneamente a mecanismos de transporte advectivo e difusão.

Aplicações desta equação surgem em diversas áreas da física e engenharia, incluindo dinâmica de fluidos, transferência de calor, dispersão de poluentes, transporte de espécies químicas e modelação atmosférica.

O principal objetivo deste trabalho consiste na resolução numérica da equação de advecção--difusão unidimensional através de métodos de diferenças finitas, analisando diferentes regimes físicos e comparando o comportamento de diferentes esquemas numéricos.

Inicialmente foi utilizado o esquema upwind para a discretização do termo advectivo. Contudo, devido à presença de difusão numérica artificial em regimes dominados pela advecção, optou-se posteriormente pela implementação do esquema de Lax--Wendroff, permitindo melhorar significativamente a precisão da solução numérica.

Além disso, foi estudada a influência relativa entre os efeitos de advecção e difusão através da análise dos parâmetros adimensionais de Courant e difusão.

\section{Fundamentos Teóricos}

\subsection{Equação de Advecção--Difusão}

A equação de advecção--difusão unidimensional pode ser escrita na forma

\[
\frac{\partial u}{\partial t}
+
v\frac{\partial u}{\partial x}
=
D\frac{\partial^2 u}{\partial x^2},
\]

onde:

\begin{itemize}
    \item $u(x,t)$ representa a quantidade transportada;
    \item $v$ corresponde à velocidade de advecção;
    \item $D$ representa o coeficiente de difusão.
\end{itemize}

O termo advectivo descreve o transporte da grandeza ao longo do domínio espacial, enquanto o termo difusivo representa o espalhamento da solução devido a efeitos dissipativos.

\subsection{Regimes Físicos}

A relação entre os efeitos de advecção e difusão pode ser analisada através dos parâmetros adimensionais

\[
C
=
\frac{v\Delta t}{\Delta x},
\]

e

\[
r
=
\frac{D\Delta t}{(\Delta x)^2}.
\]

O parâmetro $C$ corresponde ao número de Courant, associado ao transporte advectivo, enquanto $r$ representa o efeito relativo da difusão.

Quando

\[
C \gg r,
\]

o regime é dominado pela advecção.

Por outro lado, quando

\[
r \gg C,
\]

a difusão torna-se dominante.

No caso intermédio,

\[
C \approx r,
\]

os efeitos de transporte e dissipação tornam-se comparáveis.

\subsection{Perfil Inicial}

Como condição inicial foi utilizado um perfil Gaussiano da forma

\[
u(x,0)
=
\exp
\left(
-\frac{(x-x_0)^2}{\sigma^2}
\right),
\]

onde:

\begin{itemize}
    \item $x_0$ representa a posição inicial do pico;
    \item $\sigma$ controla a largura inicial da distribuição.
\end{itemize}

Este perfil permite observar simultaneamente fenómenos de deslocamento e espalhamento ao longo do tempo.

\section{Métodos Numéricos}

A resolução numérica da equação foi realizada através do método das diferenças finitas.

Considerou-se uma malha uniforme definida por

\[
x_i = i\Delta x,
\qquad
t^n = n\Delta t,
\]

sendo a solução aproximada por

\[
u_i^n \approx u(x_i,t^n).
\]

\subsection{Discretização Temporal}

A derivada temporal foi aproximada através do método de Euler explícito:

\[
\frac{\partial u}{\partial t}
\approx
\frac{u_i^{n+1}-u_i^n}{\Delta t}.
\]

Este método permite calcular iterativamente a solução no instante seguinte a partir dos valores conhecidos no instante atual.

\subsection{Discretização do Termo Difusivo}

O termo difusivo foi discretizado utilizando diferenças centradas de segunda ordem:

\[
\frac{\partial^2 u}{\partial x^2}
\approx
\frac{
u_{i+1}^n
-
2u_i^n
+
u_{i-1}^n
}{
(\Delta x)^2
}.
\]

Esta aproximação apresenta boa precisão espacial e estabilidade adequada para o problema estudado.

\subsection{Esquema Upwind}

Inicialmente foi implementado o esquema upwind para o termo advectivo.

Para velocidades positivas $(v>0)$, a derivada espacial foi aproximada por

\[
\frac{\partial u}{\partial x}
\approx
\frac{
u_i^n-u_{i-1}^n
}{
\Delta x
}.
\]

A discretização completa torna-se

\[
u_i^{n+1}
=
u_i^n
-
\frac{v\Delta t}{\Delta x}
\left(
u_i^n-u_{i-1}^n
\right)
+
\frac{D\Delta t}{(\Delta x)^2}
\left(
u_{i+1}^n-2u_i^n+u_{i-1}^n
\right).
\]

O esquema upwind apresenta elevada estabilidade numérica, especialmente em regimes dominados pela advecção. Contudo, introduz difusão numérica artificial, provocando o alisamento progressivo da solução e redução artificial da amplitude do pico Gaussiano.

\subsection{Esquema de Lax--Wendroff}

De forma a reduzir os efeitos de difusão numérica observados no esquema upwind, foi posteriormente implementado o esquema de Lax--Wendroff para a discretização do termo advectivo.

Este método utiliza uma expansão de Taylor de segunda ordem no tempo, conduzindo à seguinte aproximação:

\[
u_i^{n+1}
=
u_i^n
-
\frac{C}{2}
\left(
u_{i+1}^n-u_{i-1}^n
\right)
+
\frac{C^2}{2}
\left(
u_{i+1}^n-2u_i^n+u_{i-1}^n
\right)
+
r
\left(
u_{i+1}^n-2u_i^n+u_{i-1}^n
\right).
\]

O esquema de Lax--Wendroff apresenta maior precisão espacial relativamente ao esquema upwind, preservando melhor a forma e amplitude do perfil transportado.

Contudo, em regiões de elevado gradiente espacial, este método pode gerar pequenas oscilações numéricas.

\subsection{Condições de Estabilidade}

A estabilidade da solução depende da escolha adequada dos parâmetros numéricos.

Para o esquema explícito utilizado, é necessário garantir que os parâmetros adimensionais permaneçam dentro de intervalos estáveis.

O número de Courant é dado por

\[
C
=
\frac{v\Delta t}{\Delta x},
\]

enquanto o parâmetro difusivo é definido por

\[
r
=
\frac{D\Delta t}{(\Delta x)^2}.
\]

A escolha apropriada destes parâmetros revelou-se fundamental para evitar instabilidades e garantir resultados fisicamente coerentes.

\section{Resultados Numéricos}

Nesta secção são apresentados os resultados obtidos para diferentes regimes físicos da equação de advecção--difusão.

As simulações foram realizadas considerando três casos distintos:

\begin{itemize}
    \item regime dominado pela advecção;
    \item regime dominado pela difusão;
    \item regime equilibrado.
\end{itemize}

\subsection{Regime Dominado pela Advecção}

No regime advectivo verificou-se um deslocamento significativo do perfil Gaussiano ao longo do domínio espacial.

A utilização do esquema upwind revelou inicialmente uma redução artificial da amplitude máxima do perfil, consequência da difusão numérica introduzida pelo método.

Após a implementação do esquema de Lax--Wendroff observou-se uma preservação significativamente superior da forma e amplitude da solução.

\begin{figure}[H]
    \centering
    \includegraphics[width=0.8\textwidth]{Gráficos/advecção_dominante.png}
    \caption{Evolução temporal da solução num regime dominado pela advecção.}
\end{figure}

\subsection{Regime Dominado pela Difusão}

No regime difusivo observou-se um espalhamento progressivo da solução, acompanhado da diminuição da amplitude máxima.

O perfil apresentou reduzido deslocamento espacial, evidenciando o predomínio do mecanismo dissipativo relativamente ao transporte advectivo.

\begin{figure}[H]
    \centering
    \includegraphics[width=0.8\textwidth]{Gráficos/difusão_dominante.png}
    \caption{Evolução temporal da solução num regime dominado pela difusão.}
\end{figure}

\subsection{Regime Equilibrado}

No regime equilibrado verificou-se simultaneamente deslocamento e espalhamento da solução.

O perfil deslocou-se ao longo do domínio devido ao transporte advectivo, sofrendo simultaneamente alargamento provocado pela difusão.

Este regime permitiu observar claramente a competição entre os dois mecanismos físicos presentes na equação.

\begin{figure}[H]
    \centering
    \includegraphics[width=0.8\textwidth]{Gráficos/equilibrado.png}
    \caption{Evolução temporal da solução para um regime equilibrado.}
\end{figure}

\subsection{Comparação Final entre Regimes}

A comparação final evidencia claramente as diferenças físicas entre os três regimes estudados.

No regime advectivo predominou o deslocamento do perfil, enquanto no regime difusivo o principal efeito observado foi o espalhamento da solução.

O regime equilibrado apresentou simultaneamente transporte e dissipação.

\begin{figure}[H]
    \centering
    \includegraphics[width=0.8\textwidth]{Gráficos/comparacao_final.png}
    \caption{Comparação final entre os diferentes regimes físicos.}
\end{figure}

\subsection{Difusão Numérica}

Durante a implementação do esquema upwind verificou-se a presença de difusão numérica artificial.

Este fenómeno conduz a um espalhamento não físico da solução, particularmente evidente em regimes dominados pela advecção.

A substituição do esquema upwind pelo método de Lax--Wendroff permitiu reduzir significativamente este efeito, melhorando a preservação da amplitude e da forma do perfil Gaussiano.

Esta comparação evidenciou a importância da escolha adequada do método numérico na resolução de equações diferenciais parciais.

\section{Conclusão}

Neste trabalho foi estudada numericamente a equação de advecção--difusão unidimensional através de métodos de diferenças finitas.

Foram analisados diferentes regimes físicos associados aos efeitos relativos de transporte advectivo e difusão.

Inicialmente foi implementado o esquema upwind, observando-se boa estabilidade numérica, embora acompanhado de difusão numérica artificial.

Posteriormente, a implementação do esquema de Lax--Wendroff permitiu melhorar significativamente a precisão da solução, reduzindo os efeitos dissipativos artificiais observados no regime advectivo.

Os resultados obtidos permitiram compreender a influência dos parâmetros adimensionais de Courant e difusão na evolução temporal da solução.

Além disso, o trabalho possibilitou aprofundar conceitos fundamentais de física computacional e métodos numéricos, incluindo discretização de equações diferenciais parciais, estabilidade numérica e análise de erros computacionais.

Como possível desenvolvimento futuro, poderão ser explorados métodos de ordem superior ou abordagens modernas baseadas em inteligência artificial, como Physics-Informed Neural Networks (PINNs).

\end{document}